\chapter{The Knob That Can't Be Turned}

\epi{Every theory in this book has a dial somebody could reach for.\\ This one is welded shut, and the universe holds the key.}

\fig[1.0]{plate_locked_knob}{The knob that can't be turned. Once $\kappa$ is set by galaxies today it cancels from the ratio, and the key to $a_0$'s history is held by expansion surveys that never look at a rotation curve.}


\section{Where \texorpdfstring{$\kappa$}{kappa} goes to disappear}

Here's the law's scale again, now written as a function of cosmic time:
\begin{equation}
a_0(z)=\kappa\,c\sqrt{G\rhoDE(z)}.
\end{equation}
Divide by today's value:
\begin{equation}
\boxed{\;\frac{a_0(z)}{a_0(0)}=\sqrt{\frac{\rhoDE(z)}{\rhoDE(0)}}\;}
\end{equation}
The knob is gone. $\kappa$ cancels, $c$ cancels, $G$ cancels. What's left is a ratio of dark energy densities at two times, and that ratio is measured by surveys of the expansion of the universe (baryon acoustic oscillations, supernovae) that have nothing to do with how galaxies rotate. Nobody gets to set it with galaxy data. Not me, not a referee, not a rival.

That's the knob that can't be turned. Once $\kappa$ is fixed by galaxies today, the law's entire history is fixed by cosmology. If dark energy is a cosmological constant, $a_0$ is exactly the same at every redshift. If dark energy evolves, $a_0$ evolves with it, by an amount you can compute in advance.

\section{What evolving dark energy does to \texorpdfstring{$a_0$}{a0}}

\begin{mathbox}{From $w(z)$ to $a_0(z)$}
Energy conservation for a fluid with pressure $p=w\rho c^2$ in an expanding universe gives
\[
\frac{d\rho}{dt}=-3H(1+w)\rho\quad\Longrightarrow\quad\frac{\rhoDE(a)}{\rhoDE(1)}=\exp\left[3\int_a^1\frac{1+w(a')}{a'}\,da'\right],
\]
where $a=1/(1+z)$ is the scale factor. For the common parameterisation $w(a)=w_0+w_a(1-a)$, the integral is elementary:
\[
\frac{\rhoDE(z)}{\rhoDE(0)}=(1+z)^{3(1+w_0+w_a)}\exp\left(-\frac{3w_az}{1+z}\right).
\]
So the law predicts, with no galaxy input at all,
\[
\frac{a_0(z)}{a_0(0)}=(1+z)^{\frac32(1+w_0+w_a)}\exp\left(-\frac{3w_az}{2(1+z)}\right).
\]
For $w_0=-1$, $w_a=0$ this is exactly 1.
\end{mathbox}

Try an illustrative evolving case, $w_0=-0.75$ and $w_a=-0.8$, roughly the kind of values recent surveys have hinted at (these are for illustration, not a fit):
\begin{center}
\begin{tabular}{lccccc}
\toprule
$z$ & 0.5 & 1 & 1.5 & 2 & 2.33\\
\midrule
$a_0(z)/a_0(0)$ & 1.068 & 1.029 & 0.965 & 0.899 & 0.858\\
\bottomrule
\end{tabular}
\end{center}
It rises a little, then falls. Each number came from cosmology alone.

Figure~\ref{fig:w0wa} in Chapter~2 shows where this illustrative case sits in the dark energy knob plane. Every point in that plane gives its own $a_0(z)$ history, and none of them needs a galaxy knob.


\fig{cpl_family}{Twenty possible dark energy histories (grey), each giving a different $a_0(z)$ with zero galaxy knobs. The surveys pick the curve; the law has to follow it. A cosmological constant gives the flat red line. The rival rule $a_0\propto H(z)$ (blue) rises steeply.}

\section{The band}

The programme reconstructed $\rhoDE(z)$ from raw baryon-acoustic and supernova data with a stated distance ruler (the lane called CFG689), and converted it to the law's prediction. At 68 per cent confidence:
\begin{center}
\begin{tabular}{lccc}
\toprule
$z$ & 1 & 1.5 & 2.33\\
\midrule
the law, $a_0(z)/a_0(0)$ & 0.91 [0.76, 1.03] & 0.94 [0.71, 1.12] & 0.91 [0.33, 1.21]\\
rival, $a_0\propto H(z)$ & 1.76 & 2.35 & 3.49\\
\bottomrule
\end{tabular}
\end{center}
These are predictions from the measured dark energy density, not measured galaxy values. Constant dark energy isn't rejected by the same data (1.97$\sigma$).

\fig{a0z}{The prediction that can't be tuned. Red: the law's band from the measured dark energy history. Blue: the rival $a_0\propto H(z)$ from the same reconstruction. Dark bar: the tuned \LCDM\ recipe of Chapter~8 put in earlier universes. Gold: one published measurement at face value (see below).}

\begin{honestbox}[What the band rests on]
The band isn't free of the retired cold sector's input. It takes the matter density as $\rho_m(z)=\omega_m(1+z)^3$, with a cold mass taken from the CMB, and with positivity alone, no CMB and no Hubble constant, its own check fails. The law doesn't explain that matter term (Chapter~12).
\end{honestbox}

\section{Everybody else's scale climbs}

Now line up every acceleration scale from Part~II and ask what it does by $z=2.33$, about eleven billion years ago.

\fig{redshift_rise}{Every scale in the book at $z=2.33$, relative to today. Halo-based scales follow the expansion rate or the matter density and climb by factors of 1.3 to 5. The law follows the dark energy density and stays near 1. Only a recipe with hand terms typed in to imitate the law stays down with it.}

This is the payoff of Chapter~6. Halo accelerations are built from $G$, $M$ and $H$, so at fixed mass they grow roughly as $H^{4/3}$, and every \LCDM\ route the paper tried rises: 1.42 times for the Kaplinghat--Turner scale, 1.49 for the strongest swoop, 3.0 to 4.6 for the frozen fourteen-knob recipe. The law can't rise unless dark energy itself rises, and the measured dark energy history doesn't rise much.

\section{What to measure: Tully--Fisher in the past}

The cleanest observable is the baryonic Tully--Fisher relation, $M_b=v_f^4/Ga_0$. At fixed rotation speed, a change in $a_0$ shifts the baryonic mass by
\begin{equation}
\Delta\log_{10}M_b\big|_{v_f}=-\log_{10}\frac{a_0(z)}{a_0(0)}.
\end{equation}

\fig{btfr_shift}{The Tully--Fisher shift at fixed speed. The law (red, with its band) stays within a few hundredths of a dex of today's relation. The rival $a_0\propto H(z)$ needs galaxies about 0.25\,dex lighter at $z=1$ and 0.54\,dex lighter at $z=2.33$. The tuned \LCDM\ recipe sits lower still.}

At $z=2.33$ the law says about $+0.04\dex$; the rival says $-0.54\dex$, a factor of 3.5 in mass. That's not a subtle difference. The trouble is measuring baryonic masses of distant galaxies to a few hundredths of a dex, which nobody can yet do.

\section{The bad news first}

Before the hard part about measurement, the bad news, because it exists and it's in the record.

In 2026 the MUSE-DARK team (Ciocan and collaborators, paper III of their series) fitted the radial acceleration relation in 79 galaxies between redshift 0.33 and 1.44 and reported
\[
a_0(z\sim1)=2.38^{+0.12}_{-0.10}\times10^{-10}\ms,
\]
rising with redshift at $a_1=1.59\pm0.10$ in the same units per unit redshift (95 per cent intervals). Our band at $z=1$ on the canonical footing is $0.91\times9.36\times10^{-11}$, a factor of about 2.8 lower. An independent author, Alcock, fitted the survey's four bins as read from its figure and found a constant $a_0$ excluded at 5.3$\sigma$ if the bars are one standard deviation, and at 10$\sigma$ or more if they're 95 per cent intervals; on that reading, $a_0\propto E(z)$ is in tension too, at 4.2$\sigma$. The MUSE-DARK paper notes that frameworks tying $a_0$ to the cosmological constant are under the same pressure. That's this law. It was received before the programme's first deposit using these data, and I claim no priority on it.

\begin{honestbox}[Why I read it as not decisive, and what that reading rests on]
The survey's main fit leaves each galaxy's stellar mass free of its photometric estimate. The programme's own frozen run on a different sample of 109 galaxies (CFG262, not a direct $a_0$ measurement) found that the rise travels with the mass route: from $z=0.52$ to 1.20 the implied change is $+0.570\pm0.094\dex$ with fitted masses, $-0.115\pm0.489$ with photometric masses plus molecular gas, and $+0.036\pm0.138$ with photometric masses alone. The photometric routes are consistent with flat, and also with $H(z)$, so they rescue nothing. A redshift-dependent mass drift of 0.13--0.19\,dex would remove the rise; the survey puts its unmodelled molecular gas near 0.2\,dex, but nothing shows such a drift is present. Their companion paper finds no Tully--Fisher zero-point evolution at $z\sim1$ ($0.00\pm0.06\dex$). \textbf{On the survey's own masses, the band and the dark energy tie fail.} That reading stands until a per-galaxy reanalysis with independent baryon masses says otherwise. The programme tried that reanalysis (FN3): with the gas unmeasured, a planted rise of the survey's size still grades as consistent with the band, so the existing data cubes can't decide.
\end{honestbox}

\fig{mass_route}{Why the published rise is read as not decisive. The programme's frozen run on 109 galaxies finds the implied change in $a_0$ depends on how stellar masses are obtained: a large rise when they're fitted freely, nothing significant when they come from photometry. The answer travels with the mass route.}



\section{Why nothing on disk can decide}

The difficulty isn't the number of galaxies; it's the part of the mass error that every galaxy in a sample shares. Call it the common mode: a shared bias from the stellar-population model, the gas inventory or the distance scale, which doesn't average down.

\fig{common_mode_cartoon}{Why more galaxies don't help on their own. Left: when each galaxy's error is independent, the average of 60 lands on the truth. Right: when every galaxy shares the same bias, the average lands confidently in the wrong place.}

 The programme froze a decision rule (AT1, its arithmetic checked in the Lean proof assistant). To separate the law from $a_0\propto H$ at 3$\sigma$, the total error must be below $0.055\dex$ at $z=1$ and $0.084\dex$ at $z=2.33$. With 0.14\,dex of scatter per galaxy, the common mode $\sigma_{\rm cm}$ and the number of galaxies $N$ must satisfy, roughly,
\begin{equation}
\sigma_{\rm cm}^2+\frac{0.14^2}{N}<(\text{budget})^2.
\end{equation}
At $\sigma_{\rm cm}=0.05\dex$ the frozen rule needs 41 galaxies at $z=1$ and only 5 at $z=2.33$. At $0.075\dex$, only $z=2.33$ stays open.

\fig{decision_power}{How many galaxies it takes to decide, as a function of the sample's shared mass error, in a simplified form of the frozen rule. Near the budget the number explodes. Every high-redshift sample on disk has a common mode of 0.30 to 0.60\,dex, far off this chart.}

The programme applied the rule to every high-redshift sample it has on disk: 22 samples, 52 redshift cells. All 52 are BLOCKED, with common modes of 0.30 to 0.60\,dex and no matched local reference. That's not a failure of the law or of its rivals. It's a statement that the experiment hasn't been done yet.

\section{What would decide it}

One measurement, named but not scheduled: \textbf{a single clean deep-MOND galaxy at $z\approx2.5$ with $a_0$ measured to about 0.10\,dex}, with gas and stars weighed independently of the rotation curve. At that redshift the law and its rival differ by more than half a dex, and the tuned recipes by more. A handful of such galaxies, with a shared mass error under 0.084\,dex, would settle it.

\fig{mock_decider}{What a deciding measurement could look like (mock data, not a measurement). Five galaxies near $z\approx2.3$ with baryonic masses weighed to 0.14\,dex. The law (red) and its rival (blue dashed) are about half a dex apart, so a handful of clean galaxies would tell them apart.}



Meanwhile, three dated tests are on the calendar.

\fig{timeline}{Dated tests. The first three are scheduled data releases; none of them tests the dark energy tie directly. The last is the measurement that would.}

\textbf{ALMA U4\_28473, public 6 November 2026.} One galaxy at high redshift, with predictions frozen in PAPER66. It tests that galaxy's mass budget, not $a_0(z)$, and one disc excludes no model.

\textbf{GS4\_24110, late November 2026} (date from earlier metadata, to be rechecked). Every pairing of the alternatives tested gives at most 0.51$\sigma$ separation, so it can't decide.

\textbf{Gaia DR4, around December 2026.} Wide binary stars, separated by thousands of astronomical units, orbit each other with accelerations below $a_0$, inside the Galaxy's own field. The prediction is sealed: the law in QUMOND form, unscreened, gives a velocity ratio $\gamma_\infty=1.149$ (alternative 1.208) through the frozen pipeline. A Newtonian result at 2--30\,thousand AU would count against the bare law. (That last sentence is a reading, not a frozen line: the registered bands were written for the fuller framework, and an amendment registers a gated branch that also predicts Newton inside 30\,thousand AU.)

\fig{wide_binaries}{What Gaia DR4's wide binaries could show. Newton gives 1. The law, sealed in advance, gives 1.149. The earlier DR3 dry run's 1.205 is evidence for nothing; it was a test of the pipeline.}

\section{Why this is the point}

Every theory in this book could, in principle, reach for a knob when a measurement goes wrong. \LCDM\ can add extensions or retune feedback; the fried chicken recipes did exactly that, and they can do it again for a high-redshift sample. The landscape can always find another vacuum. Even my law has knobs it could reach for in other places: a gas term for clusters, a different interpolating function, a filter length.

But not here. The redshift dependence of $a_0$ contains no galaxy parameter. If clean measurements find $a_0$ at $z=2.5$ rising like the expansion rate, or like a halo, the law is dead, and there's nothing to turn to save it. If they find it flat, or following the dark energy history the surveys measure, the law will have done something no knob in this book could do: predict the past of galaxies from the expansion of the universe.

That's what a knob that can't be turned is for.
